\section{פתרון משוואת הדיפוזיה במרחב אינסופי}

בהרצאה זו נתמקד בפתרון האנליטי של משוואת הדיפוזיה בתחום אינסופי $x \in (-\infty, \infty)$ באמצעות טרנספורם פורייה. נדון במשמעות הפיזיקלית של הפתרון (פונקציית גרין), גבולות התקפות של המשוואה, ומבוא לבעיות בתחום סופי.

\subsection{חזרה: חוקי שימור ודעיכה באינסוף}
בשיעור הקודם עסקנו במשוואת הדיפוזיה החד-ממדית:
\begin{equation}
    \frac{\partial u}{\partial t} = a^2 \frac{\partial^2 u}{\partial x^2}
\end{equation}
כאשר $u(x,t)$ מייצגת צפיפות חלקיקים או צפיפות אנרגיה (טמפרטורה), ו-$a^2$ הוא מקדם הדיפוזיה (\LR{Diffusion Coefficient}).

הגדרנו שני גדלים שמורים פוטנציאליים:
\begin{enumerate}
    \item \textbf{המסה הכוללת (או האנרגיה):} $M(t) = \int_{-\infty}^{\infty} u(x,t) dx$.
    \item \textbf{מומנט ראשון (קשור למרכז המסה):} $P(t) = \int_{-\infty}^{\infty} x u(x,t) dx$. מרכז המסה הוא $x_{cm} = P/M$.
\end{enumerate}

\begin{theorembox}{תנאים לשימור מסה ומרכז מסה}
    על מנת שהמסה $M$ תישמר ($dM/dt = 0$), נדרש ששטף החלקיקים בקצוות יתאפס:
    \[ \left. \frac{\partial u}{\partial x} \right|_{-\infty}^{\infty} = 0 \]
    על מנת שמרכז המסה $P$ יישמר (במערכת ללא סחיפה), נדרשת דעיכה מהירה יותר של הזנבות. אם $M$ סופי, הרי ש-$u$ דועכת מהר יותר מ-$1/x$. כדי ש-$P$ יישמר, נדרש שגם האיברים הנובעים מאינטגרציה בחלקים יתאפסו:
    \[ [x u_x - u]_{-\infty}^{\infty} = 0 \]
    דבר זה מתקיים כאשר $u$ דועכת מהר יותר מ-$1/x^2$ באינסוף.
\end{theorembox}

\subsection{ניסוח הבעיה האינסופית וסיווגה}
אנו מעוניינים לפתור את בעיית הערך ההתחלתי (\LR{Initial Value Problem}):
\begin{equation}
    \begin{cases}
        u_t = a^2 u_{xx} & -\infty < x < \infty, \quad t > 0 \\
        u(x,0) = f(x) & \text{\LR{Initial Condition}}
    \end{cases}
\end{equation}
בניגוד למשוואת הגלים, כאן ניתנת רק הפונקציה $u$ בזמן $t=0$, ללא נגזרתה בזמן.

\begin{notebox}{סיווג המשוואה וקרקטריסטיקות}
    עבור המשוואה $u_t - a^2 u_{xx} = 0$, המקדמים הם $A=-a^2, B=0, C=0$ (עבור המשתנים $t,x$ בהתאמה לפורמליזם $Au_{xx} + Bu_{xt} + Cu_{tt}$).
    הדסקרימיננטה היא $\Delta = B^2 - 4AC = 0$, ולכן המשוואה היא \textbf{פרבולית} (\LR{Parabolic}).
    
    הקרקטריסטיקות נקבעות על ידי $\frac{dt}{dx} = 0 \implies t = \text{const}$.
    כלומר, הקרקטריסטיקות הן קווים מקבילים לציר ה-$x$. העובדה שאנו נותנים תנאי התחלה ($t=0$) על גבי קרקטריסטיקה אינה מובילה לסתירה (בניגוד למשוואת הגלים), כיוון שזו משוואה מסדר ראשון בזמן. איננו יכולים לתת תנאי קושי מלא (גם $u$ וגם $u_t$), שכן $u_t$ כבר נקבע על ידי $u_{xx}$ דרך המשוואה עצמה.
\end{notebox}

\subsection{פתרון באמצעות טרנספורם פורייה}
נשתמש בטרנספורם פורייה במרחב ($x$) כדי להפוך את המד"ח למד"ר.
נגדיר את הטרנספורם $R(k,t)$ ואת הטרנספורם ההפוך:
\begin{definitionbox}{טרנספורם פורייה}
    \[ R(k,t) = \mathcal{F}[u](k) = \frac{1}{\sqrt{2\pi}} \int_{-\infty}^{\infty} u(x,t) e^{-ikx} dx \]
    \[ u(x,t) = \mathcal{F}^{-1}[R](x) = \frac{1}{\sqrt{2\pi}} \int_{-\infty}^{\infty} R(k,t) e^{ikx} dk \]
\end{definitionbox}

נפעיל את הטרנספורם על משוואת הדיפוזיה:
\begin{enumerate}
    \item \textbf{אגף שמאל:} הנגזרת בזמן יוצאת מחוץ לאינטגרל:
    \[ \mathcal{F}[u_t] = \frac{\partial}{\partial t} R(k,t) \]
    \item \textbf{אגף ימין:} נגזרת מרחבית שקולה לכפל ב-$(ik)$:
    \[ \mathcal{F}[u_{xx}] = (ik)^2 R(k,t) = -k^2 R(k,t) \]
\end{enumerate}

קיבלנו משוואה דיפרנציאלית רגילה (עבור כל $k$) בזמן:
\begin{equation}
    \frac{dR}{dt} = -a^2 k^2 R
\end{equation}
פתרון המד"ר הוא דעיכה אקספוננציאלית:
\begin{equation}
    R(k,t) = C(k) e^{-a^2 k^2 t}
\end{equation}
כאשר $C(k)$ נקבע על ידי תנאי ההתחלה. בזמן $t=0$:
\[ R(k,0) = C(k) = \mathcal{F}[f(x)] \]
כלומר, $C(k)$ הוא טרנספורם פורייה של הפרופיל ההתחלתי.

\subsubsection{ניתוח המודים (Modes)}
הפתרון במרחב התדר מראה שכל מוד $k$ דועך בזמן. קצב הדעיכה תלוי ב-$k^2$.
\begin{itemize}
    \item $k$ גדול (אורך גל $\lambda$ קצר): דעיכה מהירה מאוד. פרטים קטנים וחדים מוחלקים מיד.
    \item $k$ קטן (אורך גל ארוך): דעיכה איטית. הצורה הכללית נשמרת לאורך זמן רב יותר.
\end{itemize}

\begin{examplebox}{דוגמה: פרופיל סינוסואידלי}
    עבור תנאי התחלה $u(x,0) = T_0 + A \cos(kx)$, הפתרון הוא:
    \[ u(x,t) = T_0 + A \cos(kx) e^{-a^2 k^2 t} \]
    רואים בבירור שהאמפליטודה של ההפרעה דועכת, וככל שהתדר המרחבי $k$ גבוה יותר, הדעיכה מהירה יותר.
\end{examplebox}

\subsection{פונקציית גרין והפתרון במרחב הממשי}
כדי למצוא את $u(x,t)$, עלינו לבצע טרנספורם הפוך לביטוי $R(k,t)$.
\[ u(x,t) = \frac{1}{\sqrt{2\pi}} \int_{-\infty}^{\infty} \underbrace{\hat{f}(k)}_{C(k)} e^{-a^2 k^2 t} e^{ikx} dk \]
נשתמש במשפט הקונבולוציה (\LR{Convolution Theorem}): מכפלה במרחב התדר היא קונבולוציה במרחב הזמן. אך נראה את החישוב הישיר כפי שהוצג בהרצאה.

נציב את $\hat{f}(k)$ במפורש:
\[ u(x,t) = \frac{1}{2\pi} \int_{-\infty}^{\infty} \left( \int_{-\infty}^{\infty} f(\xi) e^{-ik\xi} d\xi \right) e^{-a^2 k^2 t} e^{ikx} dk \]
נחליף את סדר האינטגרציה ונבודד את החלק שתלוי ב-$k$:
\begin{equation}
    u(x,t) = \int_{-\infty}^{\infty} f(\xi) \underbrace{\left[ \frac{1}{2\pi} \int_{-\infty}^{\infty} e^{-a^2 k^2 t + ik(x-\xi)} dk \right]}_{G(x-\xi, t)} d\xi
\end{equation}
האינטגרל הפנימי הוא אינטגרל גאוסיאני. נפתור אותו על ידי השלמה לריבוע. הביטוי באקספוננט:
\[ -a^2 t \left( k^2 - \frac{i(x-\xi)}{a^2 t} k \right) \]
נבצע השלמה לריבוע ונקבל אינטגרל גאוסיאני סטנדרטי $\int e^{-z^2} dz = \sqrt{\pi}$.
התוצאה של האינטגרל הפנימי היא הפונקציה הבאה:

\begin{resultbox}{הפתרון היסודי / פונקציית גרין (Heat Kernel)}
    \begin{equation}
        G(x,t) = \frac{1}{\sqrt{4\pi a^2 t}} e^{-\frac{x^2}{4a^2 t}}
    \end{equation}
    זהו הפתרון עבור תנאי התחלה של פונקציית דלתא ($f(x) = \delta(x)$).
\end{resultbox}

הפתרון הכללי הוא קונבולוציה של תנאי ההתחלה עם פונקציית גרין:
\begin{equation}
    u(x,t) = \int_{-\infty}^{\infty} f(\xi) G(x-\xi, t) d\xi = \frac{1}{\sqrt{4\pi a^2 t}} \int_{-\infty}^{\infty} f(\xi) e^{-\frac{(x-\xi)^2}{4a^2 t}} d\xi
\end{equation}

\subsection{ניתוח פיזיקלי של פונקציית גרין}
הפונקציה $G(x,t)$ היא גאוסיאן המשתנה בזמן.

\begin{center}
\begin{tikzpicture}
\begin{axis}[
    axis lines = left,
    xlabel = $x$,
    ylabel = {$G(x,t)$},
    ymin=0, ymax=1.2,
    xmin=-5, xmax=5,
    legend pos=north east,
    no markers,
    title={התפתחות הגאוסיאן בזמן}
]
\addplot [domain=-5:5, samples=100, blue, thick] {1/sqrt(2*pi*0.1)*exp(-x^2/(2*0.1))};
\addlegendentry{$t=t_1$ (קטן)}
\addplot [domain=-5:5, samples=100, red, thick] {1/sqrt(2*pi*0.5)*exp(-x^2/(2*0.5))};
\addlegendentry{$t=t_2$}
\addplot [domain=-5:5, samples=100, green!60!black, thick] {1/sqrt(2*pi*2)*exp(-x^2/(2*2))};
\addlegendentry{$t=t_3$ (גדול)}
\end{axis}
\end{tikzpicture}
\end{center}

תכונות מרכזיות:
\begin{enumerate}
    \item \textbf{רוחב:} סדר גודל של $\sigma \sim \sqrt{2 a^2 t}$. הרוחב גדל כמו $\sqrt{t}$. זהו המהלך האקראי (\LR{Random Walk}) - החלקיקים מתפשטים לאט יחסית לתנועה בליסטית (שם $x \sim t$).
    \item \textbf{גובה:} המקסימום קטן כמו $1/\sqrt{t}$.
    \item \textbf{שטח:} השטח מתחת לעקומה נשאר קבוע ושווה ל-1 (שימור מסה).
    \item \textbf{גבול $t \to 0$:} הפונקציה שואפת לפונקציית דלתא $\delta(x)$.
\end{enumerate}

\subsubsection{פרדוקס מהירות התפשטות אינסופית וגבולות המודל}
אם נציב $t=\epsilon$ (זמן קצר כרצוננו), הפונקציה $G(x,t)$ אינה מתאפסת באף נקודה במרחב (למעט באינסוף).
\begin{notebox}{פרדוקס הפיזיקליות}
    משמעות הדבר היא שלכאורה יש סיכוי למצוא חלקיק במרחק עצום מהמקור מיד לאחר השחרור, מה שמרמז על מהירות אינסופית (גבוהה ממהירות האור).
    
    \textbf{הסבר:} משוואת הדיפוזיה היא קירוב סטטיסטי (\LR{Continuum Approximation}). היא מניחה ממוצעים תרמיים ומספר רב של חלקיקים.
    \begin{itemize}
        \item התפלגות המהירויות בגז (מקסוול-בולצמן) דועכת אקספוננציאלית ($e^{-mv^2/2kT}$), כך שההסתברות למהירויות עצומות היא אפסית מעשית, אך לא מתמטית במודל הרציף.
        \item בזמנים קצרים מאוד או במרחקים גדולים מאוד (הזנבות), מספר החלקיקים קטן מאוד, והנחות היסוד של המשוואה (כמו חוק Fick ומושג הטמפרטורה) קורסות.
    \end{itemize}
\end{notebox}

\subsection{מבוא לבעיות בתחום סופי}
כאשר הבעיה מוגדרת בתחום סופי $x \in [0, L]$, תנאי השפה הופכים לקריטיים.

\begin{definitionbox}{סוגי תנאי שפה}
    \begin{enumerate}
        \item \textbf{תנאי דיריכלה (Dirichlet) / קצוות פתוחים:}
        קובעים את ערך הפונקציה בקצה.
        \[ u(0,t) = T_1, \quad u(L,t) = T_2 \]
        משמעות פיזיקלית: הצימוד לאמבט חום השומר על טמפרטורה קבועה, או מיכל פתוח שבו הצפיפות בקצה מתאפסת (חומר בורח). במקרה זה \textbf{המסה אינה נשמרת}.
        
        \item \textbf{תנאי נוימן (Neumann) / קצוות מבודדים:}
        קובעים את הנגזרת (השטף) בקצה.
        \[ \frac{\partial u}{\partial x}(0,t) = 0, \quad \frac{\partial u}{\partial x}(L,t) = 0 \]
        משמעות פיזיקלית: דפנות מבודדות תרמית (\LR{Insulated}), או מיכל סגור. במקרה זה \textbf{המסה נשמרת}.
    \end{enumerate}
\end{definitionbox}

\subsubsection{מצב עמיד ($t \to \infty$)}
עבור $t \to \infty$, הפתרון מתייצב ($u_t = 0$), והמשוואה הופכת ל-$u_{xx} = 0$, שפתרונה הוא קו ישר $u(x) = Ax + B$.
\begin{itemize}
    \item \textbf{קצוות מבודדים:} מכיוון שהשטף אפס, השיפוע $A=0$. הפונקציה תהיה קבועה. ערך הקבוע נקבע על ידי שימור המסה (ממוצע תנאי ההתחלה).
    \item \textbf{קצוות פתוחים (לאמבט אפס):} האנרגיה בורחת החוצה, ובסוף התהליך $u(x) = 0$.
\end{itemize}

